\chapter{现代控制理论绪论}
\label{ch:control-introduction}

\section{知识点}

\subsection{反馈与传递函数}

在单位负反馈系统中，设参考输入为 $r$，输出为 $y$，控制量为 $u$。控制器根据偏差 $e=r-y$ 产生 $u$，被控对象在 $u$ 的作用下产生新的输出。输出变化又会经反馈改变偏差，形成闭环。

对于单输入单输出线性定常系统，传递函数 $G(s)=Y(s)/U(s)$ 按零初始条件定义。已知对象 $G(s)$ 与控制器 $K(s)$，且闭环互联良定时，
\[
Y(s)=G(s)K(s)\bigl(R(s)-Y(s)\bigr),\qquad
\frac{Y(s)}{R(s)}=\frac{G(s)K(s)}{1+G(s)K(s)}.
\]
频率响应、根轨迹等经典方法常用来分析稳定性、稳定裕度与动态性能。典型指标有幅值裕度、相角裕度、超调量和调节时间。

\subsection{状态空间模型}

状态 $x(t_0)$ 和 $t_0$ 以后的输入一起确定系统后续的状态和输出。连续时间线性定常系统常写成
\[
\dot x=Ax+Bu,\qquad y=Cx+Du.
\]
状态方程描述内部运动，输出方程给出测量值。向量 $u,y$ 可分别包含多个输入和输出。对于时变系统，矩阵可以随时间变化。对于非线性系统，可用 $\dot x=f(x,u,t)$、$y=h(x,u,t)$ 建模。

对状态方程作拉普拉斯变换，整理出
\[
Y(s)=C(sI-A)^{-1}x(0)
     +\bigl[C(sI-A)^{-1}B+D\bigr]U(s).
\]
第一项是由初始状态产生的响应，方括号是系统的传递函数矩阵。给定输入输出传递函数时，状态坐标可有多种选取，系统的内部模态还要结合能控性、能观性判断。

\subsection{本课程的分析与设计问题}

\begin{itemize}
    \item 给定微分方程、结构图或传递函数，如何建立状态空间模型？选定模型以后，如何求给定初值和输入下的状态响应？
    \item 输入能否驱动每个状态方向？给定输入、输出记录，能否恢复初始状态？这分别对应能控性和能观性。
    \item 在小扰动下，状态是否保持在平衡点附近，是否随时间趋近平衡点？可用特征根或李雅普诺夫方法判断。
    \item 怎样选状态反馈得到所需的闭环极点？状态无法直接测量时，怎样利用输入和输出构造状态观测器？
\end{itemize}

讲义还介绍了几个研究方向的理论工具。李雅普诺夫方法用于稳定性分析，卡尔曼滤波用于状态估计，庞特里亚金极大值原理与贝尔曼动态规划用于最优控制。经典频域指标与状态空间中的内部性质可在同一闭环设计中配合使用。

\section{作业}

\ex{思考1（第1章讲义第3页）\quad 自动控制最重要的手段是什么？}{
    说明这种手段提供什么信息，并指出这种信息如何影响控制量。
}

\nt{\textbf{知识点}\quad 负反馈、偏差、闭环。}

\sol 自动控制中常用的核心手段是反馈。传感器测得当前输出，把测量信息送回控制端。以单位负反馈为例，求和点计算 $e=r-y$，控制器再由 $e$ 决定 $u$。在参考量固定时，若扰动使 $y$ 增大，偏差 $e$ 立即减小。控制器和对象的闭环设计决定接下来的调节效果。反馈把已经发生的输出变化纳入控制决策，使控制器能根据运行结果修正动作。

\ex{思考2（第1章讲义第3页）\quad 经典控制利用什么量，解决什么问题？}{
    以讲义中的单回路负反馈框图为例，说明控制器的输入以及控制设计的目标。
}

\nt{\textbf{知识点}\quad 输入输出描述、传递函数、闭环性能。}

\sol 控制器利用参考量与反馈量的偏差。若反馈环节为 $H(s)$，则 $E=R-HY$、$U=KE$、$Y=GU$。将三式联立可得
\[
\frac{Y(s)}{R(s)}=\frac{G(s)K(s)}{1+G(s)K(s)H(s)}.
\]
分析问题是给定控制器后判断闭环稳定性和性能。设计问题是选择控制器 $K(s)$，使闭环稳定，并满足稳态精度、响应速度、超调量或稳定裕度等要求。计算时先从框图确认求和点的正负号，再写闭环传递函数。
